
\begin{table}[h!]
\centering
\caption{Matriks komparasi ketiga kondisi terhadap kebutuhan sistem.
Kondisi C adalah solusi yang diusulkan.}
\label{tab:komparasi-solusi}
\begin{tabularx}{\textwidth}{|X|c|c|c|}
\hline
\textbf{Kebutuhan / Kriteria} &
\textbf{Kondisi A}\newline LLM Murni &
\textbf{Kondisi B}\newline RAG Konvensional &
\textbf{Kondisi C}\newline \textbf{GraphRAG} \\
\hline
Reduksi halusinasi (M1) & ✗ & Sebagian & \textbf{Ya} \\
\hline
Bobot sinyal kepercayaan SO (M2, F1) & ✗ & ✗ & \textbf{Ya} \\
\hline
Keterlacakan sumber / \emph{citation} (M3, F5) & ✗ & Sebagian & \textbf{Ya (eksplisit)} \\
\hline
Penalaran relasional / multi-hop (M4, F4) & ✗ & ✗ & \textbf{Ya} \\
\hline
Pembaruan tanpa \emph{re-training} (M5, F7) & ✗ & Ya & \textbf{Ya} \\
\hline
Kinerja pada domain Framework \& Library & Lemah & Sedang & \textbf{Target terbaik} \\
\hline
Konsistensi dengan temuan G1 (Bab II) & Tidak & Tidak & \textbf{Ya} \\
\hline
Konsistensi dengan temuan G2 (Bab II) & Tidak & Tidak & \textbf{Ya} \\
\hline
Konsistensi dengan temuan G3 (Bab II) & Tidak & Tidak & \textbf{Ya} \\
\hline
\textbf{Kesesuaian keseluruhan} & \textbf{Rendah} & \textbf{Sedang} & \textbf{Tinggi} \\
\hline
\end{tabularx}
\end{table}
